\section{Compute Verification}

\label{sec:verification}

\subsection{TEE-Based Verification}

The core verification relies on \TEE{} attestation:

\begin{theorem}[\TEE{} Attestation Soundness]
Under the \TEE{} security assumption (hardware not compromised), if attestation $A$ verifies, then the claimed computation was performed within the enclave with inputs/outputs matching the hashes.
\end{theorem}

\subsection{Efficient Re-Execution}

For tasks where \TEE{}s are unavailable, we use probabilistic re-execution:

\begin{algorithm}[t]
\caption{Probabilistic Verification}
\label{alg:prob-verify}
\begin{algorithmic}[1]
\STATE \textbf{Input}: Claimed output $y$, task $T$, sampling rate $p$
\STATE \textbf{Output}: Verified / Rejected
\IF{$\text{Random}() < p$}
    \STATE $y' \leftarrow \text{ReExecute}(T)$
    \IF{$y' \neq y$}
        \STATE Slash worker stake
        \RETURN Rejected
    \ENDIF
\ENDIF
\RETURN Verified
\end{algorithmic}
\end{algorithm}

\begin{theorem}[Detection Probability]
With sampling rate $p$ and $n$ claims, a cheating worker is caught with probability $1 - (1-p)^n$.
\end{theorem}

For $p = 0.1$ and $n = 100$, detection probability exceeds 99.99\%.

\subsection{Hybrid Verification}

We combine \TEE{} attestation with probabilistic re-execution:

\begin{equation}
V(A) = \begin{cases}
\text{TEE.Verify}(A) & \text{if TEE available} \\
\text{ProbVerify}(A, p) & \text{otherwise}
\end{cases}
\end{equation}

